\scanpage{71}
\mspage{31}
satisfying conditions a) c), \inserted{d)} of page 28.
Do we also have condition b) of page 28
\[
\Sigma_{ik}\cap\Sigma_{jk}=\varnothing\quad\hbox{if }i,j<k,\quad
i,j\ \hbox{incomparable ?}
\]
If we had $x\in\Sigma_{ik}\cap\Sigma_{jk}\inserted{\subset\Sigma_k}$,
then we would have two elements $\unsure{x_i}$ and $x_j$ of $\Sigma$,
over $i$ and $j$ respectively, which are $<x$
(over $k$), hence (since $\Sigma_{<x}$ is total\textsuperscript{ly} ordered)
\deleted{\unsure{We would have|we would have}} $x$ and $y$ would be comparable, hence also
$i$ and $j$, absurd.

\uline{Consequence} \unsure{The|the} structures ``equisingular
unfoldings of stratifications, indexed by $I$'',
described \unsure{\deleted{on}|on} page 28, are equivalent to the following structure :

\margin{\unsure{if|and} the connected components of the $\Sigma_i$ are of finite dimension...}

\textcircled{a} An \deleted{$\Sigma$} ordered topological space $\Sigma$,
satisfying \unsure{the|\deleted{the}} conditions a) b) c) and 1°) of page 27,
and \uline{moreover} : the condition that, \deleted{\unsure{for every}}
setting $\widetilde I=\pi_0(\Sigma)$, we have, for
$i,j\in\widetilde I$,
% Source-position: a detached dark point lies just above the tilde of this I; its role is unassigned.
$i<j$, that the proper
\inserted{bounding} immersion morphism $\Sigma_{ij}\overset{b_{ij}}{\unsure{\longrightarrow|\mathrel{-}}}\Sigma_j$ is
\uline{a monomorphism}

\textcircled{b} A strictly increasing map
(not necessarily a surjection)
\[
\widetilde I\longrightarrow I
\]
\margin{70}
